This chapter reports what happens when the toolkit is pointed at
benchmarks it does not own: a peer-reviewed verification paper's results
table, the standard reinforcement-learning environment suite, a hardware
model-checking benchmark, a family of semiconductor-fab process models
--- and, fusing the second and third, an Atari 2600 console compiled
into a digital circuit so that its games become white-box verification
targets. Each campaign produces \emph{certified} results: a concrete
policy, controller, or ranking function together with a
machine-checkable certificate that a stated quantitative claim holds.
(The youngest campaign carries at this preview only the weakest rung
of that ladder, and says so.)

The working standard is the Executive Summary's assurance stance
(\bookref{ch:exec-summary}{the Executive Summary}), applied to
experiments: a result earns its row by being an improvement on at least
one axis --- bound tightness, soundness level, or the cost of
\emph{checking} --- while regressing on none that matters, and by
carrying an artefact the reader can re-check rather than an assertion
the reader must trust. Where a comparison point exists in the
literature, it appears beside our number; where none exists, the row
says so; where our own result is weak or missing, that appears in the
same tables at the same prominence. Every number in this chapter is
re-derived mechanically from receipt files in the experiments
repository (section~\ref{sec:exp-repro}) at each preview; none is
transcribed by hand. The chapter is the working record of a moving
campaign: it reports state as of July 30, 2026, names what is in
flight as in-flight, and grows at each preview. One transition spans
this preview: the certificate vocabulary has been unified into a
single committed schema, and the campaigns below --- produced on its
predecessors --- are being re-conducted uniformly on it;
section~\ref{sec:exp-repro} states the discipline, and the tables
flip to re-conducted receipts as they land.

\section{What a certified experiment claims}
\label{sec:exp-claims}

``Certified'' is doing precise work in this chapter, and it resolves to
one of five claim species. Throughout, $\pi$ is a concrete shipped
policy (a lookup table, a switching law, a short program --- never a
description of a training run), and the certificate is a finite object
whose acceptance by an independent checker implies the stated bound.

\begin{description}
\item[U1 --- worst-case return over the initial set.] For deterministic
  dynamics with a set-valued initial condition: return $\ge B$ for
  \emph{every} $x_0 \in X_0$. No probability appears in the
  certificate. This is a stronger statement than any empirical average,
  and it is the natural claim for the classic control environments.
\item[U2 --- certified cost-to-goal.] For shortest-path-style
  objectives: a ranking argument bounding steps-to-goal, hence return
  $\ge -B_{\text{steps}}$.
\item[U3 --- exact or almost-sure probabilistic bounds.] For genuinely
  stochastic transitions: success probability or expected return
  $\ge B$, exact in $\mathbb{Q}$ for finite chains (backward-induction
  tables re-verified in exact arithmetic), supermartingale certificates
  for continuous state.
\item[U4 --- infinite-horizon averages.] Throughput and long-run
  stability claims for queueing-network models: fluid-limit drift
  certificates (in the tradition of Dai's fluid-limit criterion),
  reducing stability of a dispatch policy at a stated load to
  piecewise-linear inequalities over polyhedral cones --- exactly the
  Farkas/linear-programming rows the checker family of
  \bookref{ch:qtsl-manual}{the QTSL manual} replays.
\item[U5 --- i.i.d.\ confidence bounds.] An anytime-valid
  lower confidence bound on $\mathbb{E}[\text{return}]$ from rollouts
  of a fixed policy: a betting-martingale ledger whose validity rests
  on Ville's inequality. Model-free, and therefore applicable even
  where the environment's semantics is a compiled physics engine ---
  but the bound holds with confidence $1-\alpha$, not almost surely:
  the least-guaranteed species of the five. U5 rows are qualitatively
  weaker evidence than U1--U4 and are always displayed as their own
  class, never mixed into a table of exact rows.
\end{description}

Two structural facts about the species matter for reading what
follows. First, some species are \emph{terminal}: an exact
finite-chain certificate (U3 on a tabular environment) computes the
true optimum, and no future method can produce a better bound at that
claim class --- only a cheaper check. Others are
\emph{anytime-refinable}: cover-based and i.i.d.-tier certificates
trace a curve of bound quality against certificate size and check
time, and the experiments below measure those curves rather than
single points. Second, the claim class is part of the result: a table
row's meaning changes entirely if a $1-\alpha$ bound is read as an
almost-sure one, so the class label travels with every row.

One scope condition applies to every species, and it is stated here
once, plainly, because it frames what all of the results below do and
do not claim. The toolkit's certified policies come from
\emph{white-box} examination of the environment's structure --- the
transition tables, the exact simulator semantics, or (in
section~\ref{sec:exp-atari}) the environment's digital-circuit
transition relation, which the toolkit can actually examine.
\emph{System identification} --- recovering that structure from
observations of a black box --- is a separate problem, and one the
toolkit does not yet have tools for. Everything below is conditional
on being handed the mechanism; nothing below learns it.
